\section{Bogoliubov--de Gennes equations}
\label{sec:bdg}

The Fourier-space Hamiltonian of Sec.~\ref{subsec:fourier-hamiltonian} is now converted into algebraic tight-binding equations.  This intermediate step is useful because it shows explicitly how the microscopic model becomes the transfer recurrence used in Sec.~\ref{sec:transfer}.

\subsection{BdG sector and quasiparticle operator}
\label{subsec:bdg-sector}

For a fixed electron spin $\sigma$, the corresponding hole has spin $-\sigma$.  We use the quasiparticle operator
\begin{align}
    \hat\gamma_{EK\sigma}^{\dagger}
    &=
    \sum_{j,m,\alpha}
    \Big[
        u_{j\alpha}(m)
        \hat c_{jmK\alpha\sigma}^{\dagger}
        +
        s_\sigma v_{j\alpha}(m)
        \hat c_{jm,-K,\alpha,-\sigma}
    \Big].
    \label{eq:qp-operator}
\end{align}
The factor $s_\sigma$ in the hole component cancels the spin-dependent sign produced by the spin-singlet pairing commutator.  With this convention, the two independent sectors $(e_\uparrow,h_\downarrow)$ and $(e_\downarrow,h_\uparrow)$ have the same off-diagonal BdG entries $+\Delta$ and $+\Delta^*$.

The eigenoperator equation
\begin{equation}
    [\hat H^{(g)},\hat\gamma_{EK\sigma}^{\dagger}]
    =
    E\hat\gamma_{EK\sigma}^{\dagger}
    \label{eq:eigenoperator}
\end{equation}
gives, in compact one-particle notation,
\begin{equation}
    \begin{pmatrix}
        H_\sigma(K) & \Delta\\
        \Delta^\dagger & -H^*_{-\sigma}(-K)
    \end{pmatrix}
    \begin{pmatrix}
        u\\v
    \end{pmatrix}
    =
    E
    \begin{pmatrix}
        u\\v
    \end{pmatrix}.
    \label{eq:BdG-block}
\end{equation}
The lower block contains three operations: spin reversal, momentum reversal, and complex conjugation.  These operations are responsible for the signs in the hole equations below.

\subsection{Monolayer BdG equations}
\label{subsec:monolayer-bdg-equations}

For a uniform monolayer region with parameters $(\mu,h,\Delta)$ and no layer gate potential, define
\begin{equation}
    \epsilon_{e,\sigma}=E+\mu+s_\sigma h,
    \qquad
    \epsilon_{h,\sigma}=E-\mu+s_\sigma h.
    \label{eq:eps-eh}
\end{equation}
Projecting Eq.~\eqref{eq:eigenoperator} onto the $A$ and $B$ sublattices gives
\begin{align}
    \epsilon_{e,\sigma}u_A(m)
    &=
    -\eta u_B(m)-\xi u_B(m+1)+\Delta v_A(m),
    \label{eq:mono-eA}\\
    \epsilon_{e,\sigma}u_B(m)
    &=
    -\eta^*u_A(m)-\xi^*u_A(m-1)+\Delta v_B(m),
    \label{eq:mono-eB}\\
    \epsilon_{h,\sigma}v_A(m)
    &=
    +\eta v_B(m)+\xi v_B(m+1)+\Delta^*u_A(m),
    \label{eq:mono-hA}\\
    \epsilon_{h,\sigma}v_B(m)
    &=
    +\eta^*v_A(m)+\xi^*v_A(m-1)+\Delta^*u_B(m).
    \label{eq:mono-hB}
\end{align}
The electron equations inherit the minus sign of the hopping Hamiltonian, whereas the hole equations have the opposite hopping sign because they come from the block $-H^*_{-\sigma}(-K)$.  The exchange field enters both $\epsilon_{e,\sigma}$ and $\epsilon_{h,\sigma}$ with the same sign because the hole is a missing electron of spin $-\sigma$.

In the ferromagnetic lead $\Delta=0$ and $h\ne0$.  In the superconducting lead $h=0$ and $\Delta=\Delta_0e^{i\varphi}$.

\subsection{Central bilayer equations}
\label{subsec:central-bilayer-equations}

Inside the central bilayer, $h=\Delta=0$.  Define the layer-dependent scalar parameters
\begin{equation}
    x_{e,j}=E+\mu_C-V_j,
    \qquad
    x_{h,j}=\mu_C-E-V_j.
    \label{eq:central-xe-xh}
\end{equation}
The electron amplitudes satisfy
\begin{align}
    x_{e,1}u_{A1}(m)
    &=
    -\eta u_{B1}(m)-\xi u_{B1}(m+1)-t_\perp u_{B2}(m),
    \label{eq:bilayer-eA1}\\
    x_{e,1}u_{B1}(m)
    &=
    -\eta^*u_{A1}(m)-\xi^*u_{A1}(m-1),
    \label{eq:bilayer-eB1}\\
    x_{e,2}u_{A2}(m)
    &=
    -\eta u_{B2}(m)-\xi u_{B2}(m+1),
    \label{eq:bilayer-eA2}\\
    x_{e,2}u_{B2}(m)
    &=
    -\eta^*u_{A2}(m)-\xi^*u_{A2}(m-1)-t_\perp u_{A1}(m).
    \label{eq:bilayer-eB2}
\end{align}
The interlayer hopping appears only in the $A_1$ and $B_2$ equations, reflecting the AB dimer bond shown in Fig.~\ref{figgeometry}(c).

The central hole equations have the same algebraic structure after the replacement
\begin{equation}
    u\mapsto v,
    \qquad
    x_{e,j}\mapsto x_{h,j}.
    \label{eq:central-hole-map}
\end{equation}
This is the reason the central BdG transfer matrix in the next section is a direct sum of two normal bilayer transfer matrices evaluated at the two scalar arguments $\mu_C+E$ and $\mu_C-E$.
